\chapter{Conception et Présentation de l'Application}

\section{Introduction}
Ce premier chapitre présente l'application avant d'entrer dans le détail de sa réalisation. Il décrit
d'abord le plan général sur lequel elle est construite, c'est-à-dire le modèle physique de données qui
fixe les objets manipulés et leurs relations, puis le diagramme de cas d'utilisation qui relie les
acteurs aux fonctionnalités. La dernière section expose l'architecture fonctionnelle : sept rôles, un
domaine d'action distinct pour chacun et des contrôles appliqués à chaque appel.

\section{Plan de l'Application}
L'application est organisée autour d'un principe simple : tout ce qui est saisi appartient à un élève,
à une séance ou à une structure de l'établissement. Le plan qui suit montre comment ce principe se
traduit dans la base de données, puis dans les cas d'utilisation proposés aux acteurs.

\subsection{Diagramme Modèle Physique de Données}
Le modèle physique de données (\ac{MPD}) représente l'implémentation du \ac{MLD} dans le \ac{SGBD}
PostgreSQL. La base \texttt{eNote} compte 111 tables, une vue, 157 clés étrangères, 76 séquences,
34 déclencheurs et 51 fonctions. Ces objets se répartissent en six domaines fonctionnels, qui
correspondent aux six groupes de pages de l'application.

\begin{longtable}{p{4.2cm} p{10.6cm}}
\caption{Répartition des tables de la base \texttt{eNote} par domaine fonctionnel}\label{tab:domaines}\\
\toprule
\textbf{Domaine} & \textbf{Tables représentatives et rôle} \\
\midrule
\endfirsthead
\toprule
\textbf{Domaine} & \textbf{Tables représentatives et rôle} \\
\midrule
\endhead
\bottomrule
\endfoot
Identité et comptes &
\texttt{person}, \path{person_has_role}, \texttt{user}, \texttt{role}, \path{secretary_detail},
\path{student_detail}, \path{student_parent_info}, \path{student_emergency_contact},
\texttt{address}. La personne physique est distincte du compte qui permet de
se connecter : un même individu peut détenir plusieurs comptes, et ses rôles scolaires sont séparés du
rôle du compte. \\
Référentiels &
\texttt{country}, \texttt{wilaya}, \texttt{daira}, \texttt{commune}, \path{military_grade},
\path{academic_grade}, \texttt{profession}, \texttt{hobby}, \path{health_status},
\texttt{relation}, \path{recruitment_source}. Chaque libellé porte ses variantes anglaise, arabe et
française, et les noms de personnes portent une forme latine et une forme arabe. \\
Structure pédagogique &
\texttt{major}, \path{major_category}, \texttt{section}, \path{teaching_section},
\path{teaching_group}, \texttt{promotion}, \texttt{classroom}, \path{classroom_category},
\path{study_year}, \path{study_semester}, \path{study_week}, \path{study_day},
\path{study_session_type}, \path{study_session_time_slot}. C'est le squelette sur lequel
s'appuient l'emploi du temps, les volumes horaires et les examens. \\
Enseignement et évaluation &
\texttt{chapter}, \texttt{lesson}, \texttt{document}, \path{document_type},
\path{library_category}, \path{library_item}, \path{library_access}, \path{study_session},
\texttt{examination}, \path{examination_type}, \path{examination_group},
\path{examination_mark_history}, \path{major_semester_factor},
\path{major_hourly_volume}. Les fichiers pédagogiques sont référencés par une clé relative et non
par un chemin absolu, afin que la même base serve tous les postes. \\
Vie scolaire &
\path{student_gets_absent}, \path{absence_justification}, \path{observation_motif},
\path{student_undergoes_observation_by_teacher}, \path{secretary_punishes_student},
\path{secretary_rewards_student}, \path{punishment_type}, \path{reward_type},
\path{punishment_justification}, \path{reward_justification}, \path{director_comment}. Chaque
ligne porte sa date, son auteur et, lorsque c'est pertinent, la séance au cours de laquelle elle a été
constatée. \\
Support et exploitation &
\path{bug_report}, \texttt{changelog}, \texttt{feature}, \path{beta_test_report},
\path{beta_test_feature}, \texttt{message}, \path{message_recipient},
\path{direct_conversation}, \path{direct_message}, \path{app_notification},
\path{app_option}, \path{app_locked_page}, \path{audit_log},
\path{user_teaching_}\allowbreak\path{category_scope}. Ce domaine porte la messagerie interne, les notifications,
les préférences globales et le journal d'activité. \\
\end{longtable}

Le modèle est représenté par la figure \ref{fig:mpd}, qui situe les six domaines et la totalité des
tables, puis par les figures \ref{fig:mpd_identite} et \ref{fig:mpd_scolarite}, qui détaillent les
entités centrales, leurs clés et leurs cardinalités. Trois choix de conception méritent d'être
soulignés. Premièrement, les libellés de référence ne sont pas dupliqués dans les tables métier : une
section pointe vers une catégorie d'enseignement, qui porte ses trois traductions, ce qui évite qu'une
correction de libellé devienne une mise à jour de milliers de lignes. Deuxièmement, les noms de
personnes suivent un schéma à deux variantes -- une translittération latine et la forme arabe -- car un
nom propre n'a pas de traduction française. Troisièmement, une colonne générée conserve l'alias
historique de chaque libellé traduit, si bien que les requêtes, les rapports et les tris existants
continuent de fonctionner pendant que l'interface bascule vers la variante de la langue active.

\begin{figure}[H]
    \centering
    \resizebox{\textwidth}{!}{% Carte des domaines du MPD : 111 tables, 6 domaines.
\begin{tikzpicture}[font=\scriptsize]
\tikzset{dom/.style={draw=black!55, rounded corners=2pt, fill=blue!3, inner sep=5pt, align=left, text width=7.4cm}}
\node[dom, anchor=north west] at (0,0) {{\scriptsize\textbf{Identité et comptes}}\\[1pt]
{\tiny 21 tables, 34 clés étrangères}\\[3pt]
{\tiny\begin{tabular}{@{}>{\raggedright\arraybackslash}p{3.42cm}>{\raggedright\arraybackslash}p{3.42cm}@{}}
{\tiny person} & {\tiny direction}\\
{\tiny person\_\allowbreak{}has\_\allowbreak{}role} & {\tiny army}\\
{\tiny user} & {\tiny academic\_\allowbreak{}level}\\
{\tiny role} & {\tiny academic\_\allowbreak{}grade}\\
{\tiny secretary\_\allowbreak{}detail} & {\tiny military\_\allowbreak{}grade}\\
{\tiny student\_\allowbreak{}detail} & {\tiny educational\_\allowbreak{}certificate}\\
{\tiny student\_\allowbreak{}parent\_\allowbreak{}info} & {\tiny student\_\allowbreak{}military\_\allowbreak{}certificate}\\
{\tiny student\_\allowbreak{}emergency\_\allowbreak{}contact} & {\tiny student\_\allowbreak{}educational\_\allowbreak{}certificate}\\
{\tiny address} & {\tiny student\_\allowbreak{}combat\_\allowbreak{}outfit}\\
{\tiny department} & {\tiny student\_\allowbreak{}parade\_\allowbreak{}uniform}\\
{\tiny } & {\tiny student\_\allowbreak{}hobby}\\
\end{tabular}}};
\node[dom, anchor=north west] at (8.1,0) {{\scriptsize\textbf{Référentiels}}\\[1pt]
{\tiny 16 tables, 3 clés étrangères}\\[3pt]
{\tiny\begin{tabular}{@{}>{\raggedright\arraybackslash}p{3.42cm}>{\raggedright\arraybackslash}p{3.42cm}@{}}
{\tiny country} & {\tiny recruitment\_\allowbreak{}source}\\
{\tiny wilaya} & {\tiny student\_\allowbreak{}speciality}\\
{\tiny daira} & {\tiny punishment\_\allowbreak{}type}\\
{\tiny commune} & {\tiny reward\_\allowbreak{}type}\\
{\tiny profession} & {\tiny punishment\_\allowbreak{}justification}\\
{\tiny hobby} & {\tiny reward\_\allowbreak{}justification}\\
{\tiny health\_\allowbreak{}status} & {\tiny observation\_\allowbreak{}motif}\\
{\tiny relation} & {\tiny document\_\allowbreak{}type}\\
\end{tabular}}};
\node[dom, anchor=north west] at (0,-4.21) {{\scriptsize\textbf{Structure pédagogique}}\\[1pt]
{\tiny 31 tables, 50 clés étrangères}\\[3pt]
{\tiny\begin{tabular}{@{}>{\raggedright\arraybackslash}p{3.42cm}>{\raggedright\arraybackslash}p{3.42cm}@{}}
{\tiny study\_\allowbreak{}year} & {\tiny teaching\_\allowbreak{}group}\\
{\tiny study\_\allowbreak{}semester} & {\tiny teaching\_\allowbreak{}category}\\
{\tiny study\_\allowbreak{}week} & {\tiny promotion}\\
{\tiny study\_\allowbreak{}day} & {\tiny classroom}\\
{\tiny major} & {\tiny classroom\_\allowbreak{}category}\\
{\tiny major\_\allowbreak{}category} & {\tiny study\_\allowbreak{}session\_\allowbreak{}type}\\
{\tiny major\_\allowbreak{}belongs\_\allowbreak{}to\_\allowbreak{}major\_\allowbreak{}category} & {\tiny study\_\allowbreak{}session\_\allowbreak{}time\_\allowbreak{}slot}\\
{\tiny major\_\allowbreak{}is\_\allowbreak{}prerequisite\_\allowbreak{}of\_\allowbreak{}major} & {\tiny person\_\allowbreak{}belongs\_\allowbreak{}to\_\allowbreak{}section}\\
{\tiny major\_\allowbreak{}semester\_\allowbreak{}factor} & {\tiny person\_\allowbreak{}studies\_\allowbreak{}major}\\
{\tiny major\_\allowbreak{}hourly\_\allowbreak{}volume} & {\tiny person\_\allowbreak{}teaches\_\allowbreak{}major}\\
{\tiny major\_\allowbreak{}hourly\_\allowbreak{}volume\_\allowbreak{}by\_\allowbreak{}session\_\allowbreak{}type} & {\tiny student\_\allowbreak{}section}\\
{\tiny major\_\allowbreak{}teacher} & {\tiny student\_\allowbreak{}section\_\allowbreak{}history}\\
{\tiny section} & {\tiny student\_\allowbreak{}dynamic\_\allowbreak{}group\_\allowbreak{}by\_\allowbreak{}major}\\
{\tiny teaching\_\allowbreak{}section} & {\tiny user\_\allowbreak{}teaching\_\allowbreak{}category\_\allowbreak{}scope}\\
{\tiny teaching\_\allowbreak{}section\_\allowbreak{}hourly\_\allowbreak{}volume} & {\tiny teacher\_\allowbreak{}constraint}\\
{\tiny teacher\_\allowbreak{}constraint\_\allowbreak{}option} & {\tiny }\\
\end{tabular}}};
\node[dom, anchor=north west] at (8.1,-4.21) {{\scriptsize\textbf{Enseignement et évaluation}}\\[1pt]
{\tiny 17 tables, 24 clés étrangères}\\[3pt]
{\tiny\begin{tabular}{@{}>{\raggedright\arraybackslash}p{3.42cm}>{\raggedright\arraybackslash}p{3.42cm}@{}}
{\tiny chapter} & {\tiny examination}\\
{\tiny lesson} & {\tiny examination\_\allowbreak{}type}\\
{\tiny document} & {\tiny examination\_\allowbreak{}group}\\
{\tiny lesson\_\allowbreak{}has\_\allowbreak{}document} & {\tiny examination\_\allowbreak{}mark\_\allowbreak{}history}\\
{\tiny document\_\allowbreak{}access} & {\tiny person\_\allowbreak{}has\_\allowbreak{}examination}\\
{\tiny study\_\allowbreak{}session} & {\tiny person\_\allowbreak{}is\_\allowbreak{}guard\_\allowbreak{}in\_\allowbreak{}examination}\\
{\tiny study\_\allowbreak{}session\_\allowbreak{}register} & {\tiny classroom\_\allowbreak{}seat}\\
{\tiny library\_\allowbreak{}category} & {\tiny library\_\allowbreak{}item}\\
{\tiny library\_\allowbreak{}access} & {\tiny }\\
\end{tabular}}};
\node[dom, anchor=north west] at (0,-9.46) {{\scriptsize\textbf{Vie scolaire}}\\[1pt]
{\tiny 6 tables, 17 clés étrangères}\\[3pt]
{\tiny\begin{tabular}{@{}>{\raggedright\arraybackslash}p{3.42cm}>{\raggedright\arraybackslash}p{3.42cm}@{}}
{\tiny student\_\allowbreak{}gets\_\allowbreak{}absent} & {\tiny secretary\_\allowbreak{}punishes\_\allowbreak{}student}\\
{\tiny absence\_\allowbreak{}justification} & {\tiny secretary\_\allowbreak{}rewards\_\allowbreak{}student}\\
{\tiny student\_\allowbreak{}undergoes\_\allowbreak{}observation\_\allowbreak{}by\_\allowbreak{}teacher} & {\tiny director\_\allowbreak{}comment}\\
\end{tabular}}};
\node[dom, anchor=north west] at (8.1,-9.46) {{\scriptsize\textbf{Support et exploitation}}\\[1pt]
{\tiny 20 tables, 29 clés étrangères}\\[3pt]
{\tiny\begin{tabular}{@{}>{\raggedright\arraybackslash}p{3.42cm}>{\raggedright\arraybackslash}p{3.42cm}@{}}
{\tiny bug\_\allowbreak{}report} & {\tiny direct\_\allowbreak{}message}\\
{\tiny changelog} & {\tiny app\_\allowbreak{}notification}\\
{\tiny feature} & {\tiny app\_\allowbreak{}notification\_\allowbreak{}read}\\
{\tiny beta\_\allowbreak{}test\_\allowbreak{}report} & {\tiny director\_\allowbreak{}notification\_\allowbreak{}read}\\
{\tiny beta\_\allowbreak{}test\_\allowbreak{}feature} & {\tiny user\_\allowbreak{}notification\_\allowbreak{}mute}\\
{\tiny message} & {\tiny app\_\allowbreak{}option}\\
{\tiny message\_\allowbreak{}recipient} & {\tiny app\_\allowbreak{}locked\_\allowbreak{}page}\\
{\tiny message\_\allowbreak{}rule} & {\tiny audit\_\allowbreak{}log}\\
{\tiny message\_\allowbreak{}target\_\allowbreak{}role} & {\tiny user\_\allowbreak{}seen\_\allowbreak{}changelog}\\
{\tiny direct\_\allowbreak{}conversation} & {\tiny app\_\allowbreak{}locked\_\allowbreak{}page\_\allowbreak{}role}\\
\end{tabular}}};
\node[anchor=north west, font=\scriptsize, text width=14.6cm, align=left] at (0,-13.46) {%
Les flèches des figures suivantes vont de la clé étrangère vers la table référencée. La lettre N marque le côté multiple ; les tables intermédiaires portent une clé primaire composée. Les trois tables de la bibliothèque (\path{library_category}, \path{library_item}, \path{library_access}) sont créées à la première ouverture de la page et ne portent pas de clé étrangère.};
\end{tikzpicture}
}
    \caption{Carte des domaines du modèle physique de données : 111 tables réparties en six domaines}
    \label{fig:mpd}
\end{figure}

\begin{figure}[H]
    \centering
    \resizebox{\textwidth}{!}{% MPD détaillé, généré depuis le schéma PostgreSQL.
\begin{tikzpicture}[font=\scriptsize]
\tikzset{ent/.style={draw=black!65, rounded corners=2pt, fill=white, inner sep=2pt, align=left},
  lien/.style={-Latex, draw=black!55, line width=0.4pt}}
\node[ent, anchor=north west] (user) at (0,0) {%
\begin{tabular}{@{}>{\raggedright\arraybackslash}p{2.85cm}>{\raggedright\arraybackslash}p{0.9cm}@{}}
\rowcolor{blue!8}\multicolumn{2}{@{}l}{\tiny\textbf{user}}\\[1pt]
{\tiny PK user\_\allowbreak{}id} & {\tiny int}\\
{\tiny person\_\allowbreak{}id} & {\tiny bigint}\\
{\tiny user\_\allowbreak{}username} & {\tiny varchar}\\
{\tiny user\_\allowbreak{}password\_\allowbreak{}hash} & {\tiny varchar}\\
{\tiny user\_\allowbreak{}email} & {\tiny varchar}\\
{\tiny is\_\allowbreak{}active} & {\tiny bool}\\
{\tiny created\_\allowbreak{}at} & {\tiny ts}\\
\end{tabular}};
\node[ent, anchor=north west] (person) at (4.4,0) {%
\begin{tabular}{@{}>{\raggedright\arraybackslash}p{2.85cm}>{\raggedright\arraybackslash}p{0.9cm}@{}}
\rowcolor{blue!8}\multicolumn{2}{@{}l}{\tiny\textbf{person}}\\[1pt]
{\tiny PK person\_\allowbreak{}id} & {\tiny bigint}\\
{\tiny FK academic\_\allowbreak{}grade\_\allowbreak{}id} & {\tiny int}\\
{\tiny FK person\_\allowbreak{}grade\_\allowbreak{}id} & {\tiny int}\\
{\tiny person\_\allowbreak{}birthdate} & {\tiny date}\\
{\tiny person\_\allowbreak{}sex} & {\tiny varchar}\\
{\tiny military\_\allowbreak{}grade\_\allowbreak{}id} & {\tiny int}\\
{\tiny person\_\allowbreak{}photo} & {\tiny varchar}\\
\end{tabular}};
\node[ent, anchor=north west] (person_has_role) at (8.8,0) {%
\begin{tabular}{@{}>{\raggedright\arraybackslash}p{2.85cm}>{\raggedright\arraybackslash}p{0.9cm}@{}}
\rowcolor{blue!8}\multicolumn{2}{@{}l}{\tiny\textbf{person\_\allowbreak{}has\_\allowbreak{}role}}\\[1pt]
{\tiny PK/FK person\_\allowbreak{}id} & {\tiny bigint}\\
{\tiny PK/FK role\_\allowbreak{}id} & {\tiny int}\\
{\tiny PK role\_\allowbreak{}start\_\allowbreak{}date} & {\tiny date}\\
{\tiny FK department\_\allowbreak{}id} & {\tiny int}\\
{\tiny role\_\allowbreak{}end\_\allowbreak{}date} & {\tiny date}\\
\end{tabular}};
\node[ent, anchor=north west] (role) at (13.2,0) {%
\begin{tabular}{@{}>{\raggedright\arraybackslash}p{2.85cm}>{\raggedright\arraybackslash}p{0.9cm}@{}}
\rowcolor{blue!8}\multicolumn{2}{@{}l}{\tiny\textbf{role}}\\[1pt]
{\tiny PK role\_\allowbreak{}id} & {\tiny int}\\
{\tiny role\_\allowbreak{}code} & {\tiny varchar}\\
{\tiny role\_\allowbreak{}label\_\allowbreak{}en} & {\tiny varchar}\\
{\tiny role\_\allowbreak{}label\_\allowbreak{}ar} & {\tiny varchar}\\
\end{tabular}};
\node[ent, anchor=north west] (student_detail) at (4.4,-4.1) {%
\begin{tabular}{@{}>{\raggedright\arraybackslash}p{2.85cm}>{\raggedright\arraybackslash}p{0.9cm}@{}}
\rowcolor{blue!8}\multicolumn{2}{@{}l}{\tiny\textbf{student\_\allowbreak{}detail}}\\[1pt]
{\tiny PK/FK person\_\allowbreak{}id} & {\tiny bigint}\\
{\tiny FK section\_\allowbreak{}id} & {\tiny int}\\
{\tiny FK student\_\allowbreak{}speciality\_\allowbreak{}id} & {\tiny int}\\
{\tiny FK student\_\allowbreak{}academic\_\allowbreak{}level\_\allowbreak{}id} & {\tiny int}\\
{\tiny FK student\_\allowbreak{}health\_\allowbreak{}status\_\allowbreak{}id} & {\tiny int}\\
{\tiny FK student\_\allowbreak{}army\_\allowbreak{}id} & {\tiny int}\\
{\tiny FK student\_\allowbreak{}birth\_\allowbreak{}place\_\allowbreak{}id} & {\tiny int}\\
{\tiny FK student\_\allowbreak{}personal\_\allowbreak{}address\_\allowbreak{}id} & {\tiny int}\\
\end{tabular}};
\node[ent, anchor=north west] (section) at (8.8,-4.1) {%
\begin{tabular}{@{}>{\raggedright\arraybackslash}p{2.85cm}>{\raggedright\arraybackslash}p{0.9cm}@{}}
\rowcolor{blue!8}\multicolumn{2}{@{}l}{\tiny\textbf{section}}\\[1pt]
{\tiny PK section\_\allowbreak{}id} & {\tiny int}\\
{\tiny FK promotion\_\allowbreak{}id} & {\tiny int}\\
{\tiny FK teaching\_\allowbreak{}category\_\allowbreak{}id} & {\tiny int}\\
{\tiny section\_\allowbreak{}code} & {\tiny varchar}\\
{\tiny section\_\allowbreak{}label\_\allowbreak{}en} & {\tiny varchar}\\
{\tiny section\_\allowbreak{}label\_\allowbreak{}ar} & {\tiny varchar}\\
\end{tabular}};
\draw[lien] (person.east) -- node[above, font=\tiny, inner sep=1pt] {N} (person_has_role.west);
\draw[lien] (person_has_role.east) -- node[above, font=\tiny, inner sep=1pt] {N} (role.west);
\draw[lien, dashed] (person.west) -- node[above, font=\tiny, inner sep=1pt] {0,1} node[below, font=\tiny, inner sep=1pt] {lien applicatif} (user.east);
\draw[lien] (person.south) -- node[right, font=\tiny, inner sep=1pt] {N} (student_detail.north);
\draw[lien] (student_detail.east) -- node[above, font=\tiny, inner sep=1pt] {N} (section.west);
\end{tikzpicture}
}
    \caption{Modèle physique de données des entités d'identité et de comptes}
    \label{fig:mpd_identite}
\end{figure}

\begin{figure}[H]
    \centering
    \resizebox{\textwidth}{!}{% MPD détaillé, généré depuis le schéma PostgreSQL.
\begin{tikzpicture}[font=\scriptsize]
\tikzset{ent/.style={draw=black!65, rounded corners=2pt, fill=white, inner sep=2pt, align=left},
  lien/.style={-Latex, draw=black!55, line width=0.4pt}}
\node[ent, anchor=north west] (major) at (0,0) {%
\begin{tabular}{@{}>{\raggedright\arraybackslash}p{2.85cm}>{\raggedright\arraybackslash}p{0.9cm}@{}}
\rowcolor{blue!8}\multicolumn{2}{@{}l}{\tiny\textbf{major}}\\[1pt]
{\tiny PK major\_\allowbreak{}id} & {\tiny int}\\
{\tiny FK major\_\allowbreak{}category\_\allowbreak{}id} & {\tiny int}\\
{\tiny FK major\_\allowbreak{}teacher\_\allowbreak{}id} & {\tiny bigint}\\
{\tiny major\_\allowbreak{}code} & {\tiny varchar}\\
{\tiny major\_\allowbreak{}label\_\allowbreak{}en} & {\tiny varchar}\\
{\tiny major\_\allowbreak{}label\_\allowbreak{}ar} & {\tiny varchar}\\
\end{tabular}};
\node[ent, anchor=north west] (teaching_group) at (4.8,0) {%
\begin{tabular}{@{}>{\raggedright\arraybackslash}p{2.85cm}>{\raggedright\arraybackslash}p{0.9cm}@{}}
\rowcolor{blue!8}\multicolumn{2}{@{}l}{\tiny\textbf{teaching\_\allowbreak{}group}}\\[1pt]
{\tiny PK teaching\_\allowbreak{}group\_\allowbreak{}id} & {\tiny int}\\
{\tiny FK promotion\_\allowbreak{}id} & {\tiny int}\\
{\tiny FK major\_\allowbreak{}id} & {\tiny int}\\
{\tiny FK study\_\allowbreak{}semester\_\allowbreak{}id} & {\tiny int}\\
{\tiny FK teaching\_\allowbreak{}category\_\allowbreak{}id} & {\tiny int}\\
{\tiny teaching\_\allowbreak{}group\_\allowbreak{}code} & {\tiny varchar}\\
{\tiny teaching\_\allowbreak{}group\_\allowbreak{}label\_\allowbreak{}en} & {\tiny varchar}\\
\end{tabular}};
\node[ent, anchor=north west] (study_session) at (9.6,0) {%
\begin{tabular}{@{}>{\raggedright\arraybackslash}p{2.85cm}>{\raggedright\arraybackslash}p{0.9cm}@{}}
\rowcolor{blue!8}\multicolumn{2}{@{}l}{\tiny\textbf{study\_\allowbreak{}session}}\\[1pt]
{\tiny PK study\_\allowbreak{}session\_\allowbreak{}id} & {\tiny int}\\
{\tiny FK classroom\_\allowbreak{}id} & {\tiny int}\\
{\tiny FK study\_\allowbreak{}session\_\allowbreak{}type\_\allowbreak{}id} & {\tiny int}\\
{\tiny FK major\_\allowbreak{}id} & {\tiny int}\\
{\tiny FK teaching\_\allowbreak{}group\_\allowbreak{}id} & {\tiny int}\\
{\tiny FK instructor\_\allowbreak{}id} & {\tiny bigint}\\
{\tiny FK study\_\allowbreak{}day\_\allowbreak{}id} & {\tiny int}\\
{\tiny study\_\allowbreak{}session\_\allowbreak{}start\_\allowbreak{}time} & {\tiny time}\\
\end{tabular}};
\node[ent, anchor=north west] (examination) at (0,-4.1) {%
\begin{tabular}{@{}>{\raggedright\arraybackslash}p{2.85cm}>{\raggedright\arraybackslash}p{0.9cm}@{}}
\rowcolor{blue!8}\multicolumn{2}{@{}l}{\tiny\textbf{examination}}\\[1pt]
{\tiny PK examination\_\allowbreak{}id} & {\tiny int}\\
{\tiny FK major\_\allowbreak{}id} & {\tiny int}\\
{\tiny FK study\_\allowbreak{}semester\_\allowbreak{}id} & {\tiny int}\\
{\tiny FK examination\_\allowbreak{}type\_\allowbreak{}id} & {\tiny int}\\
{\tiny examination\_\allowbreak{}code} & {\tiny varchar}\\
{\tiny examination\_\allowbreak{}start\_\allowbreak{}datetime} & {\tiny ts}\\
{\tiny examination\_\allowbreak{}end\_\allowbreak{}datetime} & {\tiny ts}\\
\end{tabular}};
\node[ent, anchor=north west] (absence_justification) at (4.8,-4.1) {%
\begin{tabular}{@{}>{\raggedright\arraybackslash}p{2.85cm}>{\raggedright\arraybackslash}p{0.9cm}@{}}
\rowcolor{blue!8}\multicolumn{2}{@{}l}{\tiny\textbf{absence\_\allowbreak{}justification}}\\[1pt]
{\tiny PK absence\_\allowbreak{}justification\_\allowbreak{}id} & {\tiny int}\\
{\tiny absence\_\allowbreak{}justification\_\allowbreak{}code} & {\tiny varchar}\\
{\tiny absence\_\allowbreak{}justification\_\allowbreak{}label\_\allowbreak{}en} & {\tiny varchar}\\
{\tiny absence\_\allowbreak{}justification\_\allowbreak{}label\_\allowbreak{}ar} & {\tiny varchar}\\
\end{tabular}};
\node[ent, anchor=north west] (student_gets_absent) at (9.6,-4.1) {%
\begin{tabular}{@{}>{\raggedright\arraybackslash}p{2.85cm}>{\raggedright\arraybackslash}p{0.9cm}@{}}
\rowcolor{blue!8}\multicolumn{2}{@{}l}{\tiny\textbf{student\_\allowbreak{}gets\_\allowbreak{}absent}}\\[1pt]
{\tiny PK/FK person\_\allowbreak{}id} & {\tiny bigint}\\
{\tiny PK/FK study\_\allowbreak{}session\_\allowbreak{}id} & {\tiny int}\\
{\tiny FK absence\_\allowbreak{}justification\_\allowbreak{}id} & {\tiny int}\\
{\tiny FK per\_\allowbreak{}person\_\allowbreak{}id} & {\tiny bigint}\\
{\tiny teacher\_\allowbreak{}absence\_\allowbreak{}note} & {\tiny text}\\
\end{tabular}};
\draw[lien] (major.east) -- node[above, font=\tiny, inner sep=1pt] {N} (teaching_group.west);
\draw[lien] (teaching_group.east) -- node[above, font=\tiny, inner sep=1pt] {N} (study_session.west);
\draw[lien] (study_session.south) -- node[right, font=\tiny, inner sep=1pt] {N} (student_gets_absent.north);
\draw[lien] (student_gets_absent.west) -- node[above, font=\tiny, inner sep=1pt] {N} (absence_justification.east);
\draw[lien] (major.south) -- node[right, font=\tiny, inner sep=1pt] {N} (examination.north);
\draw[lien] (study_session.north) -- ++(0,0.55) -| node[pos=0.25, above, font=\tiny, inner sep=1pt] {N} (major.north);
\end{tikzpicture}
}
    \caption{Modèle physique de données des entités de scolarité et d'évaluation}
    \label{fig:mpd_scolarite}
\end{figure}

\subsection{Diagramme de Cas d'Utilisation}
Le diagramme de cas d'utilisation identifie les acteurs et les fonctionnalités mises à leur disposition
(figure \ref{fig:usecase}). Sept acteurs sont reconnus : l'administrateur, le directeur, le secrétaire,
l'enseignant, l'étudiant, le planificateur et le développeur. Chaque acteur dispose d'un compte unique,
et le rôle du compte détermine à la fois les pages accessibles et les opérations autorisées. Deux cas
sont communs à tous les rôles : travailler pendant une panne du serveur, sur les copies que le poste
conserve, et reprendre une saisie interrompue.

\begin{itemize}
    \item \textbf{Administrateur} : gérer les personnes et les comptes, construire le calendrier
          académique et la structure pédagogique, verrouiller des pages, administrer les mots de passe,
          consulter les statistiques et le journal d'activité ;
    \item \textbf{Directeur} : consulter les dossiers en lecture seule, suivre les séances, les
          absences, les observations et l'assiduité des enseignants, commenter les enregistrements,
          générer les rapports individuels ;
    \item \textbf{Secrétaire} : inscrire et mettre à jour les élèves des catégories qui lui sont
          attribuées, affecter les sections, enregistrer récompenses et sanctions selon ses droits ;
    \item \textbf{Enseignant} : consulter son emploi du temps, faire l'appel, consigner les
          observations, saisir les notes d'examen, publier des supports de cours ;
    \item \textbf{Étudiant} : consulter son emploi du temps, ses notes, les supports de cours et ses
          notifications ;
    \item \textbf{Planificateur} : construire l'emploi du temps des catégories qui lui sont confiées,
          suivre les volumes horaires, programmer les examens et leur disposition en salle ;
    \item \textbf{Développeur} : traiter les rapports d'anomalie, lire les rapports de test bêta,
          consulter le journal d'activité et la console de diagnostic de la base.
\end{itemize}

\begin{figure}[H]
    \centering
    \resizebox{\textwidth}{!}{% Diagramme de cas d'utilisation : deux colonnes, liens en peigne sans traverser les ellipses.
\newcommand{\acteur}[2]{%
  \begin{scope}[shift={#1}]
    \draw[line width=0.5pt] (0,0.32) circle (0.10);
    \draw[line width=0.5pt] (0,0.22) -- (0,-0.05);
    \draw[line width=0.5pt] (-0.16,0.06) -- (0.16,0.06);
    \draw[line width=0.5pt] (0,-0.05) -- (-0.12,-0.30);
    \draw[line width=0.5pt] (0,-0.05) -- (0.12,-0.30);
    \node[font=\scriptsize, align=center, text width=2.2cm] at (0,-0.58) {#2};
  \end{scope}%
}
\begin{tikzpicture}[font=\scriptsize,
  uc/.style={draw=black!55, ellipse, fill=white, align=center, text width=3.5cm, inner sep=1pt, minimum height=0.74cm},
  lien/.style={draw=black!35, line width=0.35pt}]
\draw[draw=black!45, rounded corners=3pt] (3.8,-6.05) rectangle (15.2,10.35);
\node[anchor=south, font=\scriptsize\bfseries] at (9.5,10.35) {eNote};

% colonne de gauche : administration, secretariat, enseignement, fonctions communes
\node[uc, fill=blue!10] (uc1) at (6.8,9.9) {Gérer les personnes};
\node[uc, fill=blue!10] (uc2) at (6.8,8.8) {Gérer les comptes};
\node[uc, fill=blue!10] (uc3) at (6.8,7.7) {Définir la structure};
\node[uc, fill=blue!10] (uc4) at (6.8,6.6) {Paramétrer les options};
\node[uc, fill=blue!10] (uc5) at (6.8,5.5) {Statistiques et journal};
\node[uc, fill=orange!16] (uc6) at (6.8,4.4) {Inscrire un élève};
\node[uc, fill=orange!16] (uc7) at (6.8,3.3) {Mettre à jour un dossier};
\node[uc, fill=orange!16] (uc8) at (6.8,2.2) {Récompense ou sanction};
\node[uc, fill=green!12] (uc9) at (6.8,1.1) {Faire l'appel};
\node[uc, fill=green!12] (uc10) at (6.8,0.0) {Saisir les notes};
\node[uc, fill=green!12] (uc11) at (6.8,-1.1) {Publier un support};
\node[uc, fill=black!7] (uc12) at (6.8,-2.2) {Choisir langue et apparence};
\node[uc, fill=black!7] (uc13) at (6.8,-3.3) {Signaler une anomalie};
\node[uc, fill=black!7] (uc27) at (6.8,-4.4) {Travailler pendant une panne};
\node[uc, fill=black!7] (uc28) at (6.8,-5.5) {Reprendre une saisie};",

% colonne de droite : etudiants, direction, planification, developpement
\node[uc, fill=violet!12] (uc14) at (12.2,9.9) {Emploi du temps et notes};
\node[uc, fill=violet!12] (uc15) at (12.2,8.8) {Lire les supports};
\node[uc, fill=violet!12] (uc16) at (12.2,7.7) {Gérer son profil};
\node[uc, fill=cyan!14] (uc17) at (12.2,6.6) {Consulter un dossier};
\node[uc, fill=cyan!14] (uc18) at (12.2,5.5) {Suivre les séances};
\node[uc, fill=cyan!14] (uc19) at (12.2,4.4) {Assiduité des enseignants};
\node[uc, fill=cyan!14] (uc20) at (12.2,3.3) {Générer un rapport PDF};
\node[uc, fill=brown!14] (uc21) at (12.2,2.2) {Construire l'emploi du temps};
\node[uc, fill=brown!14] (uc22) at (12.2,1.1) {Suivre les volumes horaires};
\node[uc, fill=brown!14] (uc23) at (12.2,0.0) {Programmer un examen};
\node[uc, fill=red!10] (uc24) at (12.2,-1.1) {Traiter les anomalies};
\node[uc, fill=red!10] (uc25) at (12.2,-2.2) {Exploiter les tests bêta};
\node[uc, fill=red!10] (uc26) at (12.2,-3.3) {Console de diagnostic};

% liens acteur -> cas : un peigne par groupe, a l'ecart des ellipses
\draw[lien] (1.45,7.7) -- (4.55,7.7);
\draw[lien] (4.55,0 |- uc1.west) -- (4.55,0 |- uc5.west);
\foreach \n in {uc1,uc2,uc3,uc4,uc5} { \draw[lien] (4.55,0 |- \n.west) -- (\n.west); }

\draw[lien] (1.45,3.3) -- (4.55,3.3);
\draw[lien] (4.55,0 |- uc6.west) -- (4.55,0 |- uc8.west);
\foreach \n in {uc6,uc7,uc8} { \draw[lien] (4.55,0 |- \n.west) -- (\n.west); }

\draw[lien] (1.45,0.0) -- (4.55,0.0);
\draw[lien] (4.55,0 |- uc9.west) -- (4.55,0 |- uc11.west);
\foreach \n in {uc9,uc10,uc11} { \draw[lien] (4.55,0 |- \n.west) -- (\n.west); }

\draw[lien] (1.45,-2.75) -- (4.55,-2.75);
\draw[lien] (4.55,0 |- uc12.west) -- (4.55,0 |- uc28.west);
\foreach \n in {uc12,uc13,uc27,uc28} { \draw[lien] (4.55,0 |- \n.west) -- (\n.west); }

\draw[lien] (17.15,8.8) -- (14.45,8.8);
\draw[lien] (14.45,0 |- uc14.east) -- (14.45,0 |- uc16.east);
\foreach \n in {uc14,uc15,uc16} { \draw[lien] (14.45,0 |- \n.east) -- (\n.east); }

\draw[lien] (17.15,4.95) -- (14.45,4.95);
\draw[lien] (14.45,0 |- uc17.east) -- (14.45,0 |- uc20.east);
\foreach \n in {uc17,uc18,uc19,uc20} { \draw[lien] (14.45,0 |- \n.east) -- (\n.east); }

\draw[lien] (17.15,1.1) -- (14.45,1.1);
\draw[lien] (14.45,0 |- uc21.east) -- (14.45,0 |- uc23.east);
\foreach \n in {uc21,uc22,uc23} { \draw[lien] (14.45,0 |- \n.east) -- (\n.east); }

\draw[lien] (17.15,-2.2) -- (14.45,-2.2);
\draw[lien] (14.45,0 |- uc24.east) -- (14.45,0 |- uc26.east);
\foreach \n in {uc24,uc25,uc26} { \draw[lien] (14.45,0 |- \n.east) -- (\n.east); }

% acteurs, places face a leur colonne
\coordinate (aadmin) at (1.45,7.7);
\coordinate (asec) at (1.45,3.3);
\coordinate (ateach) at (1.45,0.0);
\coordinate (aall) at (1.45,-2.75);
\coordinate (astud) at (17.15,8.8);
\coordinate (adir) at (17.15,4.95);
\coordinate (asched) at (17.15,1.1);
\coordinate (adev) at (17.15,-2.2);
\acteur{(1.45,7.7)}{Administrateur}
\acteur{(1.45,3.3)}{Secrétaire}
\acteur{(1.45,0.0)}{Enseignant}
\acteur{(1.45,-2.75)}{Utilisateur}
\acteur{(17.15,8.8)}{Étudiant}
\acteur{(17.15,4.95)}{Directeur}
\acteur{(17.15,1.1)}{Planificateur}
\acteur{(17.15,-2.2)}{Développeur}

% légende
\fill[blue!10, draw=black!40] (3.9,-6.55) rectangle ++(0.32,0.18);
\node[right, font=\scriptsize, inner sep=2pt] at (4.2,-6.46) {Admin.};
\fill[orange!16, draw=black!40] (7.1,-6.55) rectangle ++(0.32,0.18);
\node[right, font=\scriptsize, inner sep=2pt] at (7.4,-6.46) {Secrétaire};
\fill[green!12, draw=black!40] (10.0,-6.55) rectangle ++(0.32,0.18);
\node[right, font=\scriptsize, inner sep=2pt] at (10.3,-6.46) {Enseignant};
\fill[violet!12, draw=black!40] (13.1,-6.55) rectangle ++(0.32,0.18);
\node[right, font=\scriptsize, inner sep=2pt] at (13.4,-6.46) {Étudiant};
\fill[cyan!14, draw=black!40] (3.9,-7.05) rectangle ++(0.32,0.18);
\node[right, font=\scriptsize, inner sep=2pt] at (4.2,-6.96) {Directeur};
\fill[brown!14, draw=black!40] (7.1,-7.05) rectangle ++(0.32,0.18);
\node[right, font=\scriptsize, inner sep=2pt] at (7.4,-6.96) {Planif.};
\fill[red!10, draw=black!40] (10.0,-7.05) rectangle ++(0.32,0.18);
\node[right, font=\scriptsize, inner sep=2pt] at (10.3,-6.96) {Dév.};
\fill[black!7, draw=black!40] (13.1,-7.05) rectangle ++(0.32,0.18);
\node[right, font=\scriptsize, inner sep=2pt] at (13.4,-6.96) {Communs};
\end{tikzpicture}
}
    \caption{Diagramme de cas d'utilisation : sept acteurs métier, un acteur générique et vingt-huit cas}
    \label{fig:usecase}
\end{figure}

\section{Architecture Fonctionnelle}
L'application repose sur une architecture basée sur les rôles (\ac{RBAC}). Chaque utilisateur dispose
d'un ensemble de fonctionnalités adaptées à ses responsabilités. Le contrôle n'est pas seulement visuel :
chaque service du processus principal relit le rôle de l'appelant avant d'exécuter une écriture, et une
route masquée dans l'interface reste néanmoins refusée si elle est appelée directement.

\subsection{Rôle d'Administrateur}
L'administrateur est capable de :
\begin{itemize}
    \item créer, modifier et supprimer les personnes, leur attribuer des rôles scolaires, une photo et
          un département ;
    \item gérer les comptes : création, activation, désactivation, modification du nom d'utilisateur, du
          courriel et du mot de passe, attribution du rôle de compte et des catégories d'enseignement ;
    \item définir la structure académique : années, semestres, semaines, filières, catégories,
          sections, groupes, promotions, salles, créneaux et types de séance ;
    \item paramétrer les options globales : verrouillage de la saisie des notes, autorisation de
          suppression d'un élève, verrou de suppression des absences après le début d'une séance,
          pages verrouillées, matrice des messages, testeurs bêta, politique de la bibliothèque
          (dépôts, revue, téléchargement, catégorie obligatoire, taille maximale) et catégories ;
    \item consulter les statistiques, les rapports d'anomalie, les rapports de test bêta et le journal
          d'activité ;
    \item inspecter la base de données depuis la console de diagnostic, en lecture seule.
\end{itemize}

\subsection{Rôle de Directeur}
Le directeur est capable de :
\begin{itemize}
    \item consulter le répertoire des élèves et ouvrir un dossier en lecture seule ;
    \item consulter le catalogue de la bibliothèque documentaire en lecture seule ;
    \item lire, pour une semaine donnée, les séances avec leurs absences et leurs observations ;
    \item commenter une absence ou une observation, sachant que tous les directeurs partagent le même
          fil de commentaires ;
    \item suivre l'assiduité des enseignants et repérer les séances non tenues ;
    \item consulter l'historique d'un élève ou l'historique global, avec graphiques et filtres ;
    \item générer et imprimer le rapport individuel d'un élève au format \ac{PDF}.
\end{itemize}

\subsection{Rôle de Secrétaire}
Le secrétaire est capable de :
\begin{itemize}
    \item inscrire un nouvel élève et renseigner l'ensemble de son dossier : identité, état civil,
          adresse, santé, famille, contacts d'urgence, loisirs et unités d'enseignement ;
    \item mettre à jour la photo et les identifiants d'un élève ;
    \item affecter un élève à une section et consulter l'historique de ses affectations ;
    \item enregistrer une récompense ou une sanction, individuellement ou pour un groupe d'élèves ;
    \item agir uniquement dans les catégories d'enseignement qui lui ont été attribuées, et uniquement
          dans la limite des droits de conduite accordés à son compte.
\end{itemize}

\subsection{Rôle d'Enseignant}
L'enseignant est capable de :
\begin{itemize}
    \item consulter son emploi du temps hebdomadaire et la liste des élèves attendus dans une séance ;
    \item faire l'appel, marquer un élève absent avec sa justification, ou confirmer qu'une séance s'est
          tenue sans incident ;
    \item consigner une observation sur un élève présent, avec son motif et son commentaire ;
    \item saisir les notes des examens de son périmètre et consulter les moyennes calculées ;
    \item publier des supports de cours : diapositives, polycopiés, images et séances enregistrées ;
    \item déposer des livres, romans, leçons et vidéos dans la bibliothèque et suivre leur revue ;
    \item consulter le dossier d'un élève qu'il enseigne, sans pouvoir le modifier.
\end{itemize}

\subsection{Rôle d'Étudiant}
L'étudiant est capable de :
\begin{itemize}
    \item consulter l'emploi du temps de sa section pour une semaine donnée ;
    \item consulter ses notes, les moyennes par matière et sa moyenne semestrielle ;
    \item lire les supports de cours publiés par ses enseignants et visionner les vidéos ;
    \item consulter la bibliothèque documentaire et télécharger les documents autorisés ;
    \item consulter ses notifications et ses messages directs ;
    \item lire sa propre fiche de renseignements et changer sa photo.
\end{itemize}

\subsection{Rôle de Planificateur}
Le planificateur est capable de :
\begin{itemize}
    \item construire et corriger l'emploi du temps hebdomadaire des catégories qui lui sont confiées ;
    \item placer une séance dans un créneau libre en choisissant le type, l'enseignant, la salle et le
          groupe ;
    \item suivre le volume horaire cible d'un groupe et les heures déjà programmées, par type de séance
          et par section ;
    \item programmer un examen et organiser sa disposition en salle : constitution des groupes,
          répartition des places, désignation des surveillants ;
    \item consulter les notes d'un examen et les moyennes semestrielles de la filière.
\end{itemize}

\subsection{Rôle de Développeur}
Le développeur est capable de :
\begin{itemize}
    \item lire et trier les rapports d'anomalie, en changeant leur statut et leur priorité ;
    \item consulter les rapports de test bêta et les fonctionnalités couvertes par chaque version ;
    \item administrer la bibliothèque documentaire, ses catégories et sa file de revue ;
    \item rédiger les notes de version dans le journal des modifications ;
    \item consulter le journal d'activité et l'exporter au format \ac{CSV} ;
    \item ouvrir la console de diagnostic de la base et vérifier l'état de la connexion.
\end{itemize}

\section{Conclusion}
Ce chapitre a présenté le plan de l'application et son architecture fonctionnelle. La base de données
compte 111 tables réparties en six domaines, et sept rôles se partagent des périmètres distincts dont
les frontières sont vérifiées par le serveur. Le chapitre suivant décrit l'environnement et les choix
techniques qui rendent cette organisation opérationnelle sur un réseau local isolé.
